\section{Evolución de la evidencia a través de los checkpoints}

La solución final no surge de seleccionar el mínimo de una única tabla. Cada checkpoint reduce
un conjunto distinto de incertidumbres.

Checkpoint~01 reduce incertidumbre \emph{semántica}: qué significa cada fila, qué unidad tiene
cada raster, qué identificador puede usarse para unir fuentes y qué información está realmente
disponible. Checkpoint~02 reduce incertidumbre \emph{representacional}: qué resumen parcelario
se puede construir sin destruir provenance y qué familias muestran shift o sensibilidad.
Checkpoint~03 reduce incertidumbre \emph{de modelo global}: qué representaciones compactas y
qué inductive biases sobreviven nested CV. Checkpoint~03D vuelve después sobre esa misma
pregunta con mayor capacidad y demuestra que el tamaño del grid no era la única limitación.
Checkpoint~04 reduce incertidumbre
\emph{específica del conjunto objetivo}: qué significa la geometría de los 59 $X$ conocidos y
qué métodos transfieren mejor rendimiento hacia esos puntos.

Esta secuencia puede resumirse como:
\[
\text{semántica}
\rightarrow
\text{representación}
\rightarrow
\text{generalización}
\rightarrow
\text{transducción}
\rightarrow
\text{combinación global--local}
\rightarrow
\text{confirmación}
\rightarrow
\text{congelamiento}.
\]

\section{Qué información terminó siendo decisiva}

El hallazgo estadístico dominante es espacial. La asociación de distancia geográfica con
diferencia absoluta de rendimiento, $\rho=0.4882$, supera ampliamente a la distancia C1/C4 y
a la similitud temporal mensual. Moran's $I=0.6797$ confirma que el patrón no es únicamente
una correlación de pares: existe clustering espacial del rendimiento.

Esta estructura explica simultáneamente tres resultados que, vistos por separado, podrían
parecer desconectados:
\begin{enumerate}
\item los modelos globales de 03C tienen buen state RMSE pero grouped RMSE mucho peor;
\item LocalRidge mejora drásticamente en pseudo-competiciones donde existen vecinos
representativos;
\item GraphDirect conserva mucho mejor performance cuando se retienen municipios completos.
\end{enumerate}

El incumbent al cierre de Checkpoint~04 es, por tanto, local no porque los modelos locales
sean universalmente mejores, sino porque la topología observada del reto hace valiosa la
transferencia desde vecinos geográficos. Checkpoint~05 prueba de manera explícita si esa
ventaja local puede complementarse con expertos globales sin perder la disciplina de
validación.

\section{Qué información no fue decisiva}

\subsection{Más variables no fue equivalente a más señal}

B2 empirical completo, B5 all y otras representaciones masivas no dominan consistentemente a
C0/C1. En small-$n$, la dimensionalidad adicional puede aumentar varianza, diluir métricas de
distancia o requerir regularización más fuerte.

\subsection{La temporalidad mensual aislada fue débil}

Las trayectorias satelitales son valiosas para construir features de fenología y estado del
cultivo, pero la regla nearest-neighbor temporal aislada no tiene suficiente asociación con
$|y_i-y_j|$ para desplazar geografía. Esto no equivale a negar valor a remote sensing:
gran parte de C4 y de la capa agronómica procede precisamente de esas series.

\subsection{El proxy municipal contemporáneo no sustituyó una etiqueta parcelaria}

SIAP 2025 tiene cobertura completa y gran relevancia contextual, pero su resolución
administrativa es demasiado gruesa para actuar como target directo. El resultado es una
lección de multi-resolution data fusion: una fuente contemporánea sólo ayuda si su escala
espacial es compatible o si existe una corrección local suficientemente informativa.

\subsection{Compute no fue la restricción final}

Checkpoint~03D prueba esta afirmación directamente. El benchmark balanced ejecutó 230 outer
fits en 248.085 minutos y comparó CatBoost, XGBoost, LightGBM, ExtraTrees,
HistGradientBoosting y controles sobre cuatro representaciones deterministas.

El mejor grouped global pasó de aproximadamente 0.748 en E123/E13 a 0.715339 con HistGB G1.
Sin embargo, ese mismo modelo obtuvo state RMSE 0.549924. XGBoost G1 produjo
0.525847/0.733080. Variantes G2/G3 de CatBoost y ExtraTrees llegaron a state RMSE cercano o
inferior a 0.51, pero se degradaron hasta aproximadamente 0.86--0.93 bajo grouped.

El cómputo adicional mueve una parte de la frontera, pero no elimina el gap entre protocolos
ni crea observaciones independientes nuevas. Con $n=138$, más capacidad no sustituye una
representación defendible y una validación que refleje la estructura espacial.

\section{Lectura de las métricas del incumbent}

El RMSE target-matched de 0.487845 de Local04D debe interpretarse como una estimación interna bajo el
protocolo diseñado para parecerse al conjunto objetivo, no como una garantía del RMSE oculto
de FIRA. Las 16 máscaras comparten parcelas y no constituyen 16 realizaciones independientes
del proceso generador.

Por ello el reporte conserva simultáneamente:
\begin{itemize}
\item RMSE medio por split;
\item RMSE pooled;
\item peor split;
\item protocolos state-random/state/grouped;
\item estabilidad de hiperparámetros;
\item selección leave-one-split-out;
\item discrepancia entre finalistas.
\end{itemize}

La evidencia se considera fuerte cuando varias de estas vistas apuntan en la misma dirección,
no cuando una métrica aislada mejora por una cantidad marginal.

\section{Limitaciones estadísticas}

\subsection{Tamaño muestral}

La limitación fundamental es que sólo existen 138 rendimientos observados. Ningún número de
features o seeds crea observaciones independientes adicionales. Las inferencias sobre
diferencias menores a pocas milésimas deben tratarse con extrema cautela.

\subsection{Dependencia espacial}

Las parcelas no son IID. La propia razón por la que LocalRidge funciona implica que una
interpretación IID de la varianza de CV sería incorrecta. Municipality-grouped funciona como
stress test, pero también representa un dominio más difícil que el conjunto real en algunas
regiones.

\subsection{Dependencia entre pseudo-splits}

Los 16 target-matched repeats reutilizan etiquetas. El leave-one-split-out impide escoger una
regla con el mismo split que luego se califica, pero la información de parcelas se solapa
entre los otros 15 y el holdout. El resultado reduce optimismo de selección, no produce un
estimador completamente independiente.

\subsection{Resolución temporal}

Las métricas temporales de 04A usan trayectorias mensuales canónicas. BASIC/PRO contienen
capturas más finas; una investigación posterior podría usar secuencias de mayor resolución.
No se continuó porque la temporalidad mensual ya mostró una transferencia mucho menor que la
geografía y el objetivo del reto exigía priorizar hipótesis con mejor evidencia.

\subsection{SIAP}

El scope exacto se implementó correctamente a nivel de grano/ciclo/modalidad/CVEGEO, pero la
alineación del ciclo Primavera--Verano con abril--octubre es una decisión del proyecto.
Además, un rendimiento municipal puede incorporar composición de productores, manejo y
superficies distintas de las parcelas del reto.

\section{Limitaciones causales}

Ningún resultado del pipeline debe leerse como estimación causal. Si SOC$\times$CEC,
precipitación o NDVI aparecen en una representación útil, no se concluye que intervenir sobre
esa variable causaría la variación observada de rendimiento. El objetivo es predicción/
reconstrucción bajo datos observacionales.

La autocorrelación espacial también puede representar múltiples mecanismos latentes:
agroclima, manejo, variedad, calendario, estructura edáfica o procesos administrativos. El
grafo es un modelo de dependencia predictiva, no una identificación causal.

\section{Reproducibilidad y separación entre investigación y entrega}

El repositorio conserva explícitamente:
\begin{itemize}
\item configuraciones YAML por checkpoint;
\item código reusable bajo \artifact{src/geocebada/};
\item runners bajo \artifact{tools/};
\item artefactos generados bajo \artifact{reports/};
\item handoffs e historial metodológico;
\item registro de asistencia de IA;
\item un archivo final de predicciones separado de diagnósticos.
\end{itemize}

La entrega no debe construirse copiando valores desde una gráfica ni desde este PDF. Mientras
Checkpoint~05 permanezca pendiente, el artefacto canónico es el CSV de 04F con precisión
completa. Después de ejecutar 05, la fuente canónica será
\artifact{reports/checkpoint_05/final_predictions.csv}; ese archivo contendrá Local04D si
ningún challenger supera todos los gates, o el challenger confirmado si la promoción pasa.

\section{Incumbent al entrar a Checkpoint 05}

La cadena de evidencia de Checkpoints~01--04 produjo el incumbent:
\[
\boxed{
\widehat{\bm y}_{\Target}^{(04F)}
=
\operatorname{LocalRidge}_{C4,\ Geo,\ k=24,\ \alpha=30,\ p=1}
(X_{\Train},y_{\Train},X_{\Target})
}.
\]

La complejidad de los checkpoints justifica por qué esta regla relativamente simple sobrevivió
representaciones de alta dimensión, árboles, CatBoost, PLS, kernels, grafos, priors externos,
routing y blends. No obstante, Checkpoint~05 plantea una hipótesis diferente a las ya cerradas:
no buscar un mejor experto aislado, sino estimar si expertos globales y locales tienen errores
complementarios explotables mediante stacking cross-fitted.

Por tanto, LocalRidge C4 sigue siendo la salida canónica mientras 05 está pendiente. Un
challenger sólo puede reemplazarlo si mejora mean y pooled RMSE en el banco de desarrollo, si
el selector leave-one-development-split-out también mejora ambos criterios y si, después de
quedar fijado, vuelve a superar Local04D en mean y pooled RMSE sobre un banco target-matched
fresco generado $X$-only. El segundo banco no se utiliza para buscar otro challenger.

\section{Siguiente etapa operativa}

Después de resolver el único experimento final de Checkpoint~05, el siguiente bloque del proyecto debe
centrarse en:
\begin{enumerate}
\item empaquetado del archivo de entrega;
\item aplicación/interfaz para explorar parcelas, features y predicciones;
\item explicación de soporte/incertidumbre sin presentar discrepancia como intervalo
probabilístico calibrado;
\item documentación del uso de fuentes externas y de IA;
\item validaciones finales de esquema, IDs y reproducibilidad.
\end{enumerate}

La aplicación podrá mostrar las predicciones de expertos, pesos del mixture-of-experts,
soporte y discrepancia como contexto, pero el valor submission-facing debe leerse del CSV
canónico producido por Checkpoint~05. El bundle
\artifact{models/final/checkpoint05_model.joblib} se reserva para reproducción exacta y para
la futura aplicación Python; una API para parcelas arbitrariamente nuevas constituye un
contrato distinto.


\section{Estado final del modelado}

Checkpoint~05 cerró sin promoción y Checkpoint~03D cerró posteriormente la única duda amplia
de capacidad global. La salida competition-facing permanece Local04D en
\artifact{reports/checkpoint_05/final_predictions.csv}.

Existe un plan 05B, pero es deliberadamente estrecho: sólo puede responder si un peso pequeño
de los nuevos globales 03D añade señal bajo las mismas pseudo-competiciones target-matched.
No es autorización para continuar buscando indefinidamente modelos o hiperparámetros.

El proyecto conserva dos contratos distintos:
\begin{itemize}
\item \textbf{competencia fija-59}: Local04D y el CSV congelado;
\item \textbf{deployment global}: el finalista 03D XGBoost exportado y verificado en ONNX.
\end{itemize}
Confundir ambos contratos produciría una afirmación incorrecta sobre qué modelo generó la
entrega del reto.
